\section{Lösung / Implementierung}
\begin{frame}{Lösung / Implementierung}
    @Remy \& Knoll \& Hafeneder \& Tuch \& Lonthoff
    \begin{itemize}
        \item TODO: Umsetzung der einzelnen Komponenten beschreiben.
        \item TODO: Wichtige Implementierungsentscheidungen.
    \end{itemize}
\end{frame}

\begin{frame}{Lösung / Implementierung -- State-Chart}
    @ Hafeneder?
    % TODO: State-Chart einer relevanten Zustandsmaschine einfügen,
    %       z.B. VGA-Steuerung oder Audio-Sampling.
    %       Entweder als Grafik (\includegraphics{...}) oder als TikZ-Automat:
    \begin{figure}
        \centering
        \begin{tikzpicture}[->, >=Stealth, node distance=2.8cm, on grid, auto]
            \node[state, initial] (s0) {IDLE};
            \node[state] (s1) [right=of s0] {STATE\,1};
            \node[state] (s2) [right=of s1] {STATE\,2};
            \path
                (s0) edge node {bedingung} (s1)
                (s1) edge node {bedingung} (s2)
                (s2) edge [bend left=45] node {reset} (s0);
        \end{tikzpicture}
        \caption{Platzhalter -- State-Chart (TODO: durch reale FSM ersetzen).}
    \end{figure}
\end{frame}
